\documentclass[11pt,a4paper]{article}
\usepackage[utf8]{inputenc}
\usepackage[T1]{fontenc}
\usepackage{lmodern}
\usepackage{geometry}
\geometry{margin=2.5cm}
\usepackage{graphicx}
\usepackage{booktabs}
\usepackage{hyperref}
\usepackage{enumitem}
\usepackage[french]{babel}

\title{\textbf{Rapport Phase 2 : EDA \& Feature Engineering}\\PFE Dislog --- CLV, Segmentation \& Système de Recommandation}
\author{Projet PFE Dislog}
\date{\today}

\begin{document}
\maketitle

\begin{abstract}
Ce rapport résume les travaux réalisés lors de la Phase 2 du projet PFE Dislog. Cette phase couvre l'analyse exploratoire des données (EDA) sur l'entrepôt de données en schéma en étoile, la connexion à SQL Server, le profilage de toutes les tables de dimension et de faits, la construction des features RFM (Récence, Fréquence, Montant), l'analyse du comportement client, les résumés statistiques et distributions, ainsi que les visualisations. Le travail est implémenté dans le notebook Jupyter \texttt{02\_eda.ipynb} et utilise les données chargées lors de la Phase 1. Les résultats alimentent la Phase 3 (segmentation et modélisation du churn) et la Phase 4 (CLV et tableaux de bord).
\end{abstract}

\tableofcontents
\newpage

%------------------------------------------------------------------------------
\section{Introduction}
%------------------------------------------------------------------------------

La Phase 2 s'appuie sur l'entrepôt de données livré en Phase 1. L'objectif est de se connecter au schéma en étoile SQL Server, d'explorer les huit tables (six dimensions, deux faits), de dériver des features exploitables --- en particulier le \textbf{RFM} (Récence, Fréquence, Montant) pour la segmentation client --- et de produire des résumés statistiques et des visualisations qui éclairent les décisions métier et la modélisation en aval.

La \textbf{Phase 2} porte sur :
\begin{itemize}[noitemsep]
  \item La mise en place et la connexion à la base de données ; le chargement de toutes les tables en étoile dans l'environnement d'analyse
  \item L'analyse exploratoire de chaque table : forme, types, échantillons, comptages de nulls, comptages de clés distinctes, plages de dates et vérifications d'intégrité référentielle
  \item La construction des features RFM et des segments clients (ex. Champions, À Risque, Dormants)
  \item L'analyse du comportement client
  \item Les résumés statistiques et distributions
  \item Les visualisations (graphiques et diagrammes)
\end{itemize}

Le livrable principal est le notebook \texttt{notebooks/02\_eda.ipynb}, qui peut être exécuté avec les tables de faits complètes ou avec des données échantillonnées pour contrôler l'utilisation de la mémoire.

%------------------------------------------------------------------------------
\section{Mise en place \& Chargement des données}
%------------------------------------------------------------------------------

\subsection{Environnement et connexion}

Le notebook ajoute la racine du projet au \texttt{sys.path}, charge les variables d'environnement depuis \texttt{.env}, et utilise \texttt{src.config.get\_sqlserver\_connection\_string()} pour construire un moteur SQLAlchemy. Les paramètres de visualisation (matplotlib, seaborn) sont configurés pour correspondre au notebook de profilage (ex. \texttt{seaborn-v0\_8-whitegrid}, taille de figure, taille de police).

\subsection{Chargement des tables en étoile}

Les huit tables en étoile sont chargées via \texttt{pd.read\_sql} :
\begin{itemize}[noitemsep]
  \item \textbf{Dimensions} : DimDate, DimCustomer, DimSeller, DimProduct, DimPromotion, DimPaymentMethod (chargement complet).
  \item \textbf{Faits} : FactSales ($\sim$9,2M lignes) et FactInvoices ($\sim$1,4M lignes).
\end{itemize}

Pour éviter les \texttt{MemoryError} sur les machines avec une RAM limitée, une option \texttt{USE\_FULL\_FACTS} est proposée. Lorsqu'elle est à \texttt{False}, les tables de faits sont chargées avec \texttt{SELECT TOP (N)} (ex. 500K pour FactSales, 200K pour FactInvoices). Lorsqu'elle est à \texttt{True}, les tables complètes sont chargées ; cela n'est recommandé qu'avec suffisamment de mémoire libre (ex. 8 Go+ pour le notebook).

%------------------------------------------------------------------------------
\section{Analyse exploratoire de toutes les tables}
%------------------------------------------------------------------------------

Pour chaque table de dimension et de faits, le notebook réalise :
\begin{itemize}[noitemsep]
  \item \textbf{Forme et types} : Nombre de lignes et de colonnes ; types des colonnes.
  \item \textbf{Vue d'ensemble des colonnes} : Comptages et pourcentages de valeurs nulles par colonne (avec mise en forme visuelle : vert quand il n'y a pas de nulls, rouge quand il y en a).
  \item \textbf{Lignes d'exemple} : Les cinq premières lignes dans un tableau formaté.
  \item \textbf{Résumé numérique} : \texttt{describe()} pour les colonnes numériques.
  \item \textbf{Métriques clés (tables de faits uniquement)} : Comptages de valeurs distinctes pour les clés importantes (ex. SaleID distinct, CustomerKey pour FactSales ; InvoiceID distinct, CustomerKey pour FactInvoices).
\end{itemize}

La mise en forme est unifiée via deux fonctions utilitaires : \texttt{profile\_dataframe} pour les tables de dimension et \texttt{profile\_fact\_table} pour les tables de faits (avec des cartes de métriques clés). La sortie est en HTML avec un bandeau bleu, des libellés de section et un formatage de tableau cohérent.

\subsection{Plages de dates et vérification d'intégrité référentielle}

Après le profilage, le notebook vérifie :
\begin{itemize}[noitemsep]
  \item Les plages de \textbf{OrderDateKey} et \textbf{PaymentDateKey} (min/max) pour confirmer que les données sont dans les limites calendaires attendues (ex. 2024).
  \item L'\textbf{intégrité référentielle} : Toutes les valeurs de CustomerKey dans FactSales et FactInvoices sont vérifiées comme existant dans DimCustomer (ou un message clair est affiché si ce n'est pas le cas).
\end{itemize}

%------------------------------------------------------------------------------
\section{Features RFM}
%------------------------------------------------------------------------------

Le RFM (Récence, Fréquence, Montant) est un cadre standard pour évaluer la valeur et l'engagement des clients. Le notebook calcule :

\begin{itemize}[noitemsep]
  \item \textbf{Récence (Recency)} : Nombre de jours depuis la dernière commande (ou le dernier paiement) du client. Des valeurs plus faibles signifient une activité plus récente.
  \item \textbf{Fréquence (Frequency)} : Nombre de transactions (ex. commandes ou factures distinctes) par client.
  \item \textbf{Montant (Monetary)} : Montant total dépensé (ex. somme des totaux de lignes ou des montants de paiement) par client.
\end{itemize}

L'agrégation est effectuée au niveau client (CustomerKey). La date de référence pour la récence est généralement la date maximale de commande ou de paiement dans le jeu de données.

\subsection{Quartiles et segments}

Chaque dimension RFM est découpée en quartiles (1--4). Pour la Récence, le quartile 4 est attribué aux clients les plus récents et le 1 aux moins récents. Pour la Fréquence et le Montant, des quartiles plus élevés indiquent une fréquence et des dépenses plus importantes. Un \textbf{score RFM} composite (ex. chaîne de chiffres de quartile) et un \textbf{libellé de segment} sont ensuite attribués. Exemples de segments :
\begin{itemize}[noitemsep]
  \item \textbf{Champions} : Récence, fréquence et montant élevés (meilleurs clients).
  \item \textbf{Clients fidèles}, \textbf{Loyalistes potentiels}, \textbf{À Risque}, \textbf{Dormants}, etc., selon les combinaisons de quartiles R, F, M.
\end{itemize}

La logique de segmentation est implémentée de sorte que R élevé (ex. 3--4), F élevé et M élevé correspondent aux Champions ; R faible et F faible correspondent souvent à Dormants ou À Risque, selon M.

%------------------------------------------------------------------------------
\section{Analyse du comportement client}
%------------------------------------------------------------------------------

Au-delà du RFM, le notebook peut inclure :
\begin{itemize}[noitemsep]
  \item Des agrégations par client (ex. valeur moyenne de commande, nombre de produits distincts, fréquence de commande dans le temps).
  \item Des patterns temporels : ventes ou factures par mois, trimestre ou jour de la semaine.
  \item Des tableaux croisés ou pivots reliant les segments à d'autres dimensions (ex. région, catégorie de produit).
\end{itemize}

Ces analyses aident à interpréter les segments RFM et à préparer les features pour la Phase 3 (clustering, prédiction du churn).

%------------------------------------------------------------------------------
\section{Résumés statistiques et distributions}
%------------------------------------------------------------------------------

Le notebook fournit des résumés statistiques et des distributions pour les mesures clés, par exemple :
\begin{itemize}[noitemsep]
  \item Tendance centrale et dispersion (moyenne, médiane, écart-type, min, max) pour les colonnes numériques des tables de faits (ex. Quantity, LineTotalAmount, PaymentAmount).
  \item Distribution des valeurs RFM (ex. histogrammes ou boîtes à moustaches par segment).
  \item Comptages et proportions par segment ou par niveau de dimension.
\end{itemize}

Ces résumés soutiennent les contrôles de qualité des données et éclairent le choix des features et des modèles dans les phases suivantes.

%------------------------------------------------------------------------------
\section{Visualisations}
%------------------------------------------------------------------------------

Les visualisations sont produites avec matplotlib et seaborn, en accord avec le style défini au début du notebook. Les exemples incluent :
\begin{itemize}[noitemsep]
  \item Des diagrammes à barres ou histogrammes pour la taille des segments, les ventes par mois ou les produits phares.
  \item Des distributions de la Récence, de la Fréquence et du Montant (ex. histogrammes ou KDE).
  \item Des heatmaps ou diagrammes à barres groupées comparant les segments selon des dimensions (ex. région, méthode de paiement).
\end{itemize}

Les figures sont configurées avec des titres clairs, des libellés d'axes et des légendes afin que les résultats puissent être utilisés dans les rapports ou les présentations.

%------------------------------------------------------------------------------
\section{Problèmes \& Complexités rencontrés en Phase 2}
%------------------------------------------------------------------------------

Bien que la Phase 2 s'appuie sur un entrepôt stable, plusieurs défis pratiques et complexités ont été rencontrés :

\subsection{Travail avec de très grandes tables de faits}

FactSales ($\sim$9,2M lignes) et FactInvoices ($\sim$1,4M lignes) sont volumineuses pour une EDA interactive dans un notebook. Le chargement des tables complètes dans pandas sur une machine avec 12 Go de RAM est proche de la limite mémoire et peut déclencher \texttt{MemoryError} sur certaines configurations. Pour gérer cela :
\begin{itemize}[noitemsep]
  \item Le notebook expose un flag \texttt{USE\_FULL\_FACTS} qui contrôle si les faits complets ou un sous-ensemble échantillonné TOP $N$ sont chargés.
  \item L'échantillonnage préserve la structure et les distributions nécessaires pour prototyper le RFM et l'analyse comportementale, tandis que les exécutions sur données complètes peuvent être réservées aux environnements à forte mémoire et aux vérifications finales.
\end{itemize}

\subsection{Performance des agrégations et du calcul RFM}

Le calcul des features RFM nécessite de regrouper des millions de lignes par CustomerKey et d'agréger les dates et les montants. Des group-bys écrits naïvement peuvent être lents ou gourmands en mémoire. Le notebook adopte donc :
\begin{itemize}[noitemsep]
  \item Des expressions de group-by compactes qui ne conservent que les colonnes nécessaires pour R, F et M.
  \item La réutilisation des résultats intermédiaires au lieu de recalculer les agrégations lourdes plusieurs fois.
  \item L'application des coupes par quartile sur les données pré-agrégées au niveau client, et non sur les tables de faits brutes.
\end{itemize}

Ces choix de conception maintiennent un temps d'exécution raisonnable même avec les tables de faits complètes.

\subsection{Calibration des règles de segmentation RFM}

Il est facile de mal interpréter les quartiles RFM (ex. traiter un score de Récence faible comme << meilleur >> alors que l'échelle est inversée). Au cours de l'expérimentation, les règles de segmentation initiales ont été affinées de sorte que :
\begin{itemize}[noitemsep]
  \item Une Récence élevée (activité récente), une Fréquence élevée et un Montant élevé correspondent clairement aux << Champions >>.
  \item Les autres segments tels que << Fidèles >>, << À Risque >> et << Dormants >> suivent des seuils cohérents et pertinents pour le métier.
\end{itemize}

La documentation de ces règles dans le notebook garantit que les phases de modélisation ultérieures pourront réutiliser la même logique de segmentation sans ambiguïté.

%------------------------------------------------------------------------------
\section{Résultats et livrables}
%------------------------------------------------------------------------------

À la fin de la Phase 2 :
\begin{itemize}[noitemsep]
  \item Toutes les tables en étoile ont été profilées avec un format cohérent et stylisé (vue d'ensemble des colonnes, échantillons, résumé numérique ; métriques clés pour les faits).
  \item Les features RFM et les segments clients sont calculés et disponibles pour utilisation en aval.
  \item L'analyse du comportement client et les résumés statistiques documentent les patterns dans les données.
  \item Les visualisations illustrent les distributions, les segments et les KPIs clés.
\end{itemize}

Le notebook \texttt{02\_eda.ipynb} est le point unique pour les résultats de la Phase 2 ; il peut être ré-exécuté avec les données de faits complètes ou échantillonnées selon la RAM disponible.

%------------------------------------------------------------------------------
\section{Conclusion}
%------------------------------------------------------------------------------

La Phase 2 est terminée. L'entrepôt de données est exploré, les features RFM et les segments sont construits, et le comportement client est analysé avec des résumés statistiques et des visualisations. Les résultats alimentent la Phase 3 (segmentation et modélisation du churn) et la Phase 4 (prédiction de la CLV et tableaux de bord).

\end{document}
